\newpage
\pagestyle{empty}

\begin{center}
	\textbf{САНКТ-ПЕТЕРБУРГСКИЙ ПОЛИТЕХНИЧЕСКИЙ УНИВЕРСИТЕТ ПЕТРА ВЕЛИКОГО} \\
	\vspace{0.5em}
	Физико-механический институт \\
	\vspace{2em}
\end{center}

\begin{flushright}
	\begin{minipage}{0.35\textwidth} 
		\begin{flushleft} 
			УТВЕРЖДАЮ\\
			Руководитель ОП \\
			К.Н. Козлов \\
			\guillemetleft \underline{\hspace{0.5cm}} \guillemetright \underline{\hspace{3.25cm}} 2026 г. 
		\end{flushleft}
	\end{minipage}
\end{flushright}

\begin{center}
	\textbf{ЗАДАНИЕ} \\
	\textbf{на выполнение выпускной квалификационной работы магистра} \\
	студенту \uline{Добрецовой Е.В.}, гр. \uline{5040102/40101} \\
\end{center}

\vspace{1em}

\vspace{1em}

\noindent

\begin{enumerate}
	\item Тема работы: «Оптимизация эволюционным алгоритмом метода решения задачи Коши на базе метода Рунге-Кутты для осциллирующей функции»
	\item Срок сдачи студентом законченной работы: июнь 2026 г.
	\item Исходные данные по работе: данные вычислительных экспериментов, полученные самостоятельно
	\item Инструментальные средства:
		\begin{itemize}
			\item язык программирования Python
			\item среда разработки Microsoft Visual Studio Code
			\item система контроля версий Git
		\end{itemize}
	\item Ключевые источники литературы:
		\begin{enumerate}
			\item Anastassi, Z.A., Kosti, A.A., and Rufai, M.A. (2023). A Parametric Method Optimised for the Solution of the (21)-Dimensional Nonlinear Schrödinger Equation+. Mathematics 11, 609. \href{https://doi.org/10.3390/math11030609}{https://doi.org/10.3390/math11030609}
			
			\item Anastassi, Z.A. (2025). Evolutionary Optimisation of Runge–Kutta Methods for Oscillatory Problems. Mathematics 13, 2796. \\ \href{https://doi.org/10.3390/math13172796}{https://doi.org/10.3390/math13172796}
		\end{enumerate}
		
		\item Содержание работы (перечень подлежащих разработке вопросов):
		\begin{enumerate}
			\item Введение. Обоснование актуальности. Цели и задачи работы
			\item Обзор литературы. Основные теоретические сведения
			\item Постановка задачи
			\item Разработка методов
			\item Результаты
			\item Выводы
			\item Заключение
		\end{enumerate}
		
		\item Дата выдачи задания \underline{09.02.2026} \hfill
\end{enumerate}

\vspace{3em}


\noindent
\begin{tabularx}{\textwidth}{X c X}
	Научный руководитель &
	\underline{\hspace{5cm}} &
	К.\,Н.~Козлов \\
	&
	{\small (подпись)} & \\
\end{tabularx}

\vspace{1.5em}

\noindent
\hspace{\tabcolsep}Задание принял к исполнению
\vspace{0.75em}

\noindent
\begin{tabularx}{\textwidth}{X c X}
	Студент &
	\underline{\hspace{5cm}} &
	Е.\,В.~Добрецова \\
	&
	{\small (подпись)} & \\
\end{tabularx}